\medskip

%Dans un collège, 91 filles et 77 garçons participent à un club sciences.
%
%On souhaite former des groupes, de sorte que chaque groupe ait le même nombre de filles et le même nombre de garçons.
%
%\medskip

\begin{enumerate}
\item %Décomposer 91 et 77 en produit de facteurs premiers.
$91 = 7 \times 13$ ;

$77 = 7 \times 11$
\item %En déduire combien de groupes au maximum on peut former.

%\textbf{Argumenter la réponse en précisant la démarche.}
Les deux nombres sont des multiples de 7 ; on pourra constituer 7 groupes au maximum.
\item %Dans ce cas combien d’élèves y aura-t-il dans chaque groupe?
Chacun des groupes comportera 13 garçons et 11 filles.
\end{enumerate}
